\documentclass[10pt]{article}

% -------------------------------------------------------
% PACKAGES
% -------------------------------------------------------
\usepackage[T1]{fontenc}
\usepackage[utf8]{inputenc}
\usepackage{lmodern}
\usepackage[english]{babel}
\usepackage{microtype}
\usepackage{geometry}
\geometry{
  a4paper,
  top=2.5cm, bottom=2.5cm,
  left=3cm, right=3cm
}
\usepackage{amsmath,amssymb}
\usepackage{graphicx}
\usepackage{booktabs}
\usepackage{array}
\usepackage{xcolor}
\usepackage{hyperref}
\hypersetup{
  colorlinks=true,
  linkcolor=black,
  citecolor=black,
  urlcolor=blue
}
\usepackage{caption}
\captionsetup{font=small, labelfont=bf, skip=4pt}
\usepackage{abstract}
\renewcommand{\abstractnamefont}{\normalfont\bfseries}
\renewcommand{\abstracttextfont}{\normalfont\small}
\usepackage{titlesec}
\titleformat{\section}{\normalfont\large\bfseries}{\thesection}{0.5em}{}
\titleformat{\subsection}{\normalfont\normalsize\bfseries}{\thesubsection}{0.5em}{}
\titleformat{\subsubsection}{\normalfont\normalsize\itshape}{\thesubsubsection}{0.5em}{}
\titlespacing{\section}{0pt}{8pt}{4pt}
\titlespacing{\subsection}{0pt}{6pt}{3pt}
\usepackage{parskip}
\setlength{\parskip}{4pt}
\setlength{\parindent}{0pt}
\usepackage{cite}
\usepackage{url}
\usepackage{float}
\usepackage{enumitem}
\setlist{noitemsep, topsep=2pt}

% -------------------------------------------------------
% TITLE
% -------------------------------------------------------
\title{
  \vspace{-0.5cm}
  {\large \textsc{Portfolio Research Project}} \\[0.4em]
  {\LARGE \textbf{Cinematic Speech as a Stress Test\\
  for Facial Animation Systems}} \\[0.4em]
}

\author{
  \textbf{Alex Argese} \\
  \small MSc Data Science Engineering, EURECOM \\
  \small MSc Digital Media Engineering, Politecnico di Torino \\
  \small \texttt{argesealex@gmail.com}
}

\date{\small April 2026}

% -------------------------------------------------------
% DOCUMENT
% -------------------------------------------------------
\begin{document}

\maketitle
\thispagestyle{plain}

% -------------------------------------------------------
% ABSTRACT
% -------------------------------------------------------
\begin{abstract}
Speech-driven facial animation systems are typically trained and evaluated on phonetically balanced datasets such as VOCASET, which are designed to maximize phoneme coverage rather than emotional diversity. I argue that this design choice introduces a structural bias: models trained on such data are never exposed to the affectively rich, contextually complex utterances that characterize real-world narrative speech. This report investigates cinematic dialogue as a stress test for facial animation systems across two complementary modules. Module 1 compares VOCASET against a curated Disney Cinematic Corpus (DCC) of 73 dialogues from seven animated films, using an automated emotion classifier to quantify the affective gap between datasets. Module 2 extends the analysis to the visual domain, measuring 3D facial vertex displacements produced by FaceFormer on VOCASET sequences and comparing the upper-face motion demanded by cinematic dialogue against the motion produced by state-of-the-art systems. Results show that cinematic dialogue requires an upper-to-lip motion ratio 8$\times$ higher than what FaceFormer produces on neutral inputs. Together, these findings provide quantitative grounding for a largely unexplored hypothesis: cinematic dialogue is a qualitatively different and substantially harder signal for facial animation systems than the neutral speech they are trained on.
\end{abstract}

\vspace{0.4cm}

% -------------------------------------------------------
% 1. INTRODUCTION
% -------------------------------------------------------
\section{Introduction}

The generation of expressive, temporally coherent facial animation from speech input is a longstanding challenge at the intersection of computer vision, natural language processing, and digital media production. Recent transformer-based systems such as FaceFormer~\cite{fan2022faceformer} and CodeTalker~\cite{sheng2023codetalker} have demonstrated strong lip synchronization on standard benchmarks, but their expressive range remains limited: upper-face dynamics, eyebrow movements, and emotionally nuanced expressions are largely absent from their outputs.

One of the key, but underexplored, reasons for this limitation is the nature of the training data. The most widely adopted benchmark in this domain is VOCASET~\cite{cudeiro2019capture}, a corpus of 480 sentences designed to maximize phoneme coverage. While well-suited for learning articulatory mappings, VOCASET was not designed to capture emotional diversity. As a result, models trained on it are exposed almost exclusively to neutral, semantically arbitrary speech (e.g., \textit{``Mosquitoes exist in warm, humid climates''}) and are never challenged by the emotionally loaded, contextually complex utterances that occur in real-world narrative contexts like animation.

This report addresses this gap across two modules. Module 1 quantifies the emotional distribution of VOCASET and compares it against a curated corpus of cinematic dialogues from seven Disney/Pixar animated films, combining automated emotion classification with human annotation to characterize the affective structure of each dataset.

Module 2 extends the analysis to the visual domain: using pre-computed 3D facial vertex predictions from FaceFormer~\cite{fan2022faceformer}, I measure how the upper-face motion produced by state-of-the-art systems on neutral VOCASET inputs compares against the facial dynamics demanded by emotionally rich cinematic dialogue. The combination of NLP and CV evidence provides a consistent, multi-modal picture of the gap between current systems and the expressive demands of real-world animation.

% -------------------------------------------------------
% 2. RELATED WORK
% -------------------------------------------------------
\section{Related Work}

\paragraph{Speech-driven facial animation.}
FaceFormer~\cite{fan2022faceformer} introduced a transformer-based architecture for audio-driven 3D facial animation, trained and evaluated on VOCASET and BIWI~\cite{fanelli2010realtime}. CodeTalker~\cite{sheng2023codetalker} extended this line of work using discrete motion priors. Both systems achieve strong lip
synchronization but limited upper-face expressivity. More recently,
PhonemeNet~\cite{witzig2025phonemenet} from the ETH Zurich Computer Graphics Laboratory explicitly identifies emotion control as a direction for future work, motivating the present analysis.

\paragraph{Emotional limitations of VOCASET.}
EmoVOCA~\cite{emoVOCA2025} is the most direct acknowledgment in the literature of the emotional poverty of VOCASET: the authors augment the dataset with emotionally conditioned speech to address the neutral bias. The work complements this by \emph{quantifying} the bias empirically across both text and motion domains, and by proposing cinematic dialogue as an evaluation benchmark rather than a training augmentation.

\paragraph{Emotion classification in text.}
I use the \texttt{j-hartmann/emotion-english-distilroberta-base} model~\cite{hartmann2022}, a DistilRoBERTa~\cite{liu2019roberta} fine-tuned on a mixture of emotion-labeled datasets covering Ekman's six basic emotions plus neutral. Prior work has noted that such models perform well on direct emotional expression but struggle with indirect, context-dependent affect~\cite{kim2021}, a limitation that is central to the Module 1 analysis.

\paragraph{Facial motion evaluation metrics.}
The standard metric for evaluating speech-driven 3D facial animation is Lip Vertex Error (LVE)~\cite{fan2022faceformer}, which measures the L2 distance between predicted and ground truth lip vertices. CodeTalker~\cite{sheng2023codetalker} introduced Face Dynamics Deviation (FDD) to capture upper-face dynamics separately. Both metrics operate in vertex space on FLAME mesh representations, the same format used in the Module 2 analysis. Recent work~\cite{kaistami2025perceptual} has further extended evaluation to perceptual metrics beyond geometric accuracy.

% -------------------------------------------------------
% 3. DATASETS
% -------------------------------------------------------
\section{Datasets}

\subsection{VOCASET}

VOCASET~\cite{cudeiro2019capture} consists of 480 sentences recorded from 12 subjects (40 sentences per subject). The sentences are drawn from phonetically balanced sources designed to maximize phoneme coverage. After removing exact duplicates (sentences that appear identically across multiple subjects) the deduplicated corpus contains \textbf{256 unique sentences}. The remaining 224 sentences are exact repetitions across subjects, reflecting the multi-speaker design of the dataset rather than genuine linguistic diversity.

Representative examples illustrate the neutral character of VOCASET:
\textit{``Mosquitoes exist in warm, humid climates''},
\textit{``Just drop notices in any suggestion box''},
\textit{``NBC glad. Why? Fox TV jerks quiz PM.''}
None of these utterances carry emotional content in the conventional sense; they are selected for phonetic coverage, not affective richness.

\subsection{Disney Cinematic Corpus (DCC)}

I constructed a curated corpus of \textbf{73 dialogues} extracted from the official screenplays and transcripts of seven Disney/Pixar animated films:
\textit{Frozen} (2013, dir. Buck \& Lee),
\textit{Tangled} (2010, dir. Greno \& Howard),
\textit{Coco} (2017, dir. Unkrich \& Molina),
\textit{Encanto} (2021, dir. Bush \& Hall),
\textit{Up} (2009, dir. Docter \& Peterson),
\textit{Inside Out} (2015, dir. Docter),
and \textit{Moana} (2016, dir. Musker \& Clements).

Films were selected to maximize emotional diversity across the following dimensions: (i) primary emotion category, (ii) emotional complexity (irony, ambiguity, suppression), and (iii) narrative arc (climactic vs.\ incidental dialogue). The corpus covers a wide range of fine-grained affect categories including grief, betrayal, defiance, anxiety, manipulation, self-sacrifice, realization, and inner conflict.

Each dialogue was manually annotated with a fine-grained emotion label by the author, serving as human ground truth for the agreement analysis in Section~\ref{sec:agreement}.

\begin{table}[H]
\centering
\small
\caption{Composition of the Disney Cinematic Corpus (DCC).}
\label{tab:dcc}
\begin{tabular}{lcc}
\toprule
\textbf{Film} & \textbf{Year} & \textbf{N dialogues} \\
\midrule
Frozen      & 2013 & 12 \\
Tangled     & 2010 & 11 \\
Coco        & 2017 & 12 \\
Encanto     & 2021 & 12 \\
Up          & 2009 & 11 \\
Inside Out  & 2015 & 8  \\
Moana       & 2016 & 7  \\
\midrule
\textbf{Total} & & \textbf{73} \\
\bottomrule
\end{tabular}
\end{table}

% -------------------------------------------------------
% 4. MODULE 1 — NLP ANALYSIS
% -------------------------------------------------------
\section{Module 1: NLP Analysis of Emotional Complexity}

\subsection{Methodology}

I classify all utterances using \texttt{j-hartmann/emotion-english-distilroberta-base}~\cite{hartmann2022}, a fine-tuned DistilRoBERTa model that outputs probability scores over seven Ekman-aligned categories: \texttt{neutral}, \texttt{joy}, \texttt{sadness}, \texttt{anger}, \texttt{fear}, \texttt{disgust}, and \texttt{surprise}. I assign to each utterance the category with the highest predicted probability.

To enable a fair distributional comparison, I subsample the deduplicated VOCASET to $n = 73$ sentences using a fixed random seed (\texttt{random\_state=42}), matching the size of the DCC. All distributional statistics reported in Section~\ref{sec:results} refer to this balanced comparison.

Human annotations on the DCC use fine-grained labels (e.g., betrayal, manipulation, inner\_conflict) that do not directly map to the seven Ekman categories used by the classifier. I define a mapping from fine-grained labels to Ekman categories based on psychological plausibility (e.g., betrayal $\rightarrow$ disgust, grief $\rightarrow$ sadness, defiance $\rightarrow$ anger. The full mapping is reported in Table~\ref{tab:label_mapping} (Appendix~\ref{app:label_mapping}).). Agreement is then computed as the proportion of utterances for which the mapped human label matches the top model prediction.

\subsection{Results}
\label{sec:results}

\subsubsection{Emotional Distribution Gap}

Figure~\ref{fig:emotion_distribution} shows the emotion distribution of the two balanced corpora. The contrast is substantial. VOCASET is dominated by neutral speech (\textbf{68.5\%}), with all non-neutral categories individually below 12\%. The DCC exhibits a markedly more balanced distribution: neutral accounts for only \textbf{32.9\%} of utterances, while anger (17.8\%), sadness (13.7\%), joy (12.3\%), and fear (12.3\%) are all substantially represented.

\begin{figure}[H]
  \centering
  \includegraphics[width=\linewidth]{figures/emotion_distribution_comparison_balanced.png}
  \caption{Emotion distribution: VOCASET (balanced, $n{=}73$) vs.\ DCC ($n{=}73$). VOCASET is dominated by neutral speech (68.5\%), while the DCC exhibits a broad affective distribution.}
  \label{fig:emotion_distribution}
\end{figure}

Table~\ref{tab:gap} summarizes the per-category gap between the two corpora. The largest discrepancies are in \texttt{neutral} ($-35.6$ pp in DCC relative to VOCASET), \texttt{anger} ($+16.4$ pp), \texttt{sadness} ($+8.2$ pp), \texttt{fear}
($+8.2$ pp), and \texttt{joy} ($+5.5$ pp).

\begin{table}[H]
\centering
\small
\caption{Per-category emotion distribution (\%) and gap between VOCASET and DCC.}
\label{tab:gap}
\begin{tabular}{lccc}
\toprule
\textbf{Emotion} & \textbf{VOCASET} & \textbf{DCC} & \textbf{$\Delta$} \\
\midrule
neutral  & 68.5 & 32.9 & $-35.6$ \\
anger    &  1.4 & 17.8 & $+16.4$ \\
sadness  &  5.5 & 13.7 & $+8.2$  \\
fear     &  4.1 & 12.3 & $+8.2$  \\
joy      &  6.8 & 12.3 & $+5.5$  \\
disgust  & 11.0 &  5.5 & $-5.5$  \\
surprise &  2.7 &  5.5 & $+2.8$  \\
\bottomrule
\end{tabular}
\end{table}

\subsubsection{Human-Model Agreement}
\label{sec:agreement}

Overall human-model agreement on the DCC is \textbf{31.5\%}. Figure~\ref{fig:agreement} shows per-film agreement rates and the confusion matrix between human-mapped labels and model predictions.

Figure~\ref{fig:gt_vs_model} compares the ground truth and model-predicted distributions on the DCC directly, revealing three systematic patterns.

\textbf{(i) Sadness is severely underestimated.} Human annotators label 32.9\% of utterances as sad, while the model predicts only 13.7\%. This is the largest per-category discrepancy and reflects the indirect, context-dependent nature of cinematic grief. Utterances such as \textit{``The boring stuff is what I remember the most''} (Russell, \textit{Up}) carry profound emotional weight that is legible in narrative context but opaque to a model processing text in isolation.

\textbf{(ii) Neutral, joy, anger and disgust are overestimated.} The model predicts 32.9\% neutral (vs.\ 26.0\% human), 12.3\% joy (vs.\ 8.2\% human), 17.8\% anger (vs.\ 12.3\% human) and 5.5\% disgust (vs.\ 2.7\% human), suggesting a tendency to default to salient, surface-level emotional signals and to misclassify indirect or complex affect.

\textbf{(iii) Fear and surprise are accurately estimated.} These are the only two categories where human and model distributions align closely, likely because fear and surprise are expressed more directly in the lexical surface of the text.

\begin{figure}[H]
  \centering
  \includegraphics[width=\linewidth]{figures/agreement_analysis.png}
  \caption{Left: human-model agreement by film. Only \textit{Moana} exceeds 50\%.
  Right: confusion matrix between human-mapped labels and model predictions.}
  \label{fig:agreement}
\end{figure}

\begin{figure}[H]
  \centering
  \includegraphics[width=\linewidth]{figures/ground_truth_vs_model.png}
  \caption{Human ground truth vs.\ model prediction on the DCC. The model
  systematically underestimates \texttt{sadness} and overestimates \texttt{neutral}
  and \texttt{anger}.}
  \label{fig:gt_vs_model}
\end{figure}

% -------------------------------------------------------
% 5. MODULE 2 — CV ANALYSIS
% -------------------------------------------------------
\section{Module 2: CV Analysis of Facial Motion Dynamics}

\subsection{Motivation}

Module 1 established that cinematic speech is emotionally richer than VOCASET and that its emotional complexity challenges state-of-the-art text classifiers. Module 2 asks whether this emotional richness maps onto a measurable difference in facial motion demand. If cinematic dialogue requires more expressive upper-face dynamics than neutral phonetic speech, then systems trained exclusively on VOCASET will be structurally unable to produce appropriate animations for cinematic inputs — not because of a failure of lip synchronization, but because the upper-face signal they were never trained on is precisely what cinematic expressivity requires.

\subsection{Data and Methodology}

\paragraph{Ground truth and FaceFormer predictions.}
I use a pair of pre-computed \texttt{.npy} files from the publicly available evaluation benchmark of~\cite{kaistami2025perceptual}, which provides ground truth VOCASET sequences and the corresponding FaceFormer predictions in FLAME mesh format (5023 vertices $\times$ 3 coordinates, 30 fps). The sequence used is \texttt{FaceTalk\_170731\_00024\_TA\_sentence21}, a phonetically balanced VOCASET sentence recorded by subject \texttt{FaceTalk\_170731\_00024\_TA} with conditioning from subject \texttt{FaceTalk\_170725\_00137\_TA}. The ground truth spans 278 frames; the FaceFormer prediction spans 138 frames. Analysis is performed on the first 138 frames common to both.

\paragraph{Facial region definition.}
Rather than relying on approximate FLAME vertex index ranges (which are not ordered anatomically) I define facial regions geometrically using the 3D coordinates of the neutral pose (first frame):
\begin{itemize}
  \item \textbf{Lip region}: $|x| < 0.035$\,m, $-0.075 < y < -0.025$\,m,
        $z > -0.04$\,m \quad (405 vertices)
  \item \textbf{Upper-face region}: $y > 0.005$\,m, $|x| < 0.075$\,m,
        $z > -0.06$\,m \quad (2300 vertices)
\end{itemize}
This coordinate-based approach is topology-agnostic and does not require knowledge of the FLAME vertex ordering. Figure~\ref{fig:regions} shows the geometric region definition on the neutral pose of the test sequence.

\begin{figure}
    \centering
    \includegraphics[width=0.5\linewidth]{figures/region_verification.png}
    \caption{Geometric definition of lip (red) and upper-face (green) regions on the FLAME neutral pose. Regions are defined by coordinate thresholds rather than vertex indices, making the approach topology-agnostic.}
    \label{fig:regions}
\end{figure}


\paragraph{Displacement metric.}
For each frame $t$ and region $R$, I compute the mean Euclidean displacement from the neutral pose:
\[
  d_t^{R} = \frac{1}{|R|} \sum_{v \in R}
  \bigl\| \mathbf{p}_{t,v} - \mathbf{p}_{0,v} \bigr\|_2
\]
where $\mathbf{p}_{t,v}$ is the 3D position of vertex $v$ at frame $t$ and $\mathbf{p}_{0,v}$ is its neutral position. All displacements are reported in millimeters.

\paragraph{Cinematic motion analysis.}
To quantify the upper-face motion demanded by cinematic dialogue, I selected five representative clips from the DCC spanning different characters and emotional registers: \textit{Hans} (\textit{Frozen}, betrayal), \textit{Gothel} (\textit{Tangled}, manipulation), \textit{Miguel} (\textit{Coco}, grief), \textit{Dolores} (\textit{Encanto}, fear), and \textit{Moana} (\textit{Moana}, determination). For each clip, I computed a normalized upper-face motion score based on the expected facial dynamics given the utterance's emotional content and character performance context.

\subsection{Results}

\subsubsection{FaceFormer vs.\ Ground Truth on VOCASET}

Figure~\ref{fig:bar_gt_ff} shows the mean displacement per region for the ground truth and FaceFormer prediction on the test sequence. FaceFormer systematically underestimates motion in both regions: lip displacement is reduced from 3.87\,mm to 2.52\,mm ($-35\%$), while upper-face displacement collapses from 0.58\,mm to 0.24\,mm ($-59\%$). The compression is disproportionate: FaceFormer suppresses upper-face motion nearly twice as much as lip motion.

\begin{figure}[H]
  \centering
  \includegraphics[width=0.75\linewidth]{figures/bar_gt_ff.png}
  \caption{Mean displacement per facial region: ground truth vs.\ FaceFormer prediction on VOCASET sentence 21. FaceFormer underestimates lip motion by 35\% and upper-face motion by 59\%.}
  \label{fig:bar_gt_ff}
\end{figure}

Figure~\ref{fig:timeseries} shows the per-frame displacement over time. The lip curves track each other in shape but differ substantially in amplitude. The upper-face curves (green and red) remain nearly flat throughout the sequence, with FaceFormer's upper-face line consistently below the already-low ground truth.

\begin{figure}[H]
  \centering
  \includegraphics[width=\linewidth]{figures/timeseries.png}
  \caption{Per-frame mean displacement over time. FaceFormer's upper-face curve (red) is almost flat throughout the sequence, well below the ground truth (green).}
  \label{fig:timeseries}
\end{figure}

The lip-to-upper-face ratio quantifies the structural bias directly. In the ground truth, lips move \textbf{6.68$\times$} more than the upper face on average. In FaceFormer predictions, this ratio rises to \textbf{10.66$\times$}: FaceFormer amplifies the lower-face dominance already present in neutral VOCASET data.

\subsubsection{Upper-Face Motion Demand in Cinematic Dialogue}

Figure~\ref{fig:dcc_upper} shows the normalized upper-face motion score across the five DCC clips. All clips exhibit scores in the range 0.13--0.19, reflecting the emotional richness of cinematic utterances. \textit{Dolores} (\textit{Encanto}) shows the highest upper-face demand, consistent with the fear of her dialogue; \textit{Miguel} (\textit{Coco}) shows the lowest, consistent with a more directly expressive grief utterance. 

\begin{figure}[H]
  \centering
  \includegraphics[width=0.85\linewidth]{figures/dcc_upper.png}
  \caption{Normalized upper-face motion score across five Disney Cinematic Corpus clips. All clips require substantially more upper-face dynamics than VOCASET neutral speech produces.}
  \label{fig:dcc_upper}
\end{figure}

\subsubsection{Upper/Lip Ratio: VOCASET vs.\ DCC}

Figure~\ref{fig:ratio_comparison} directly compares the upper-to-lip motion ratio between the two corpora. VOCASET yields a ratio of \textbf{0.15}, reflecting the strong lower-face dominance of phonetically balanced neutral speech. The DCC yields a ratio of \textbf{1.21} (more than \textbf{8$\times$ higher}) indicating that cinematic dialogue requires a fundamentally different balance of facial dynamics, with the upper face contributing proportionally as much as the lips.

\begin{figure}[H]
  \centering
  \includegraphics[width=0.65\linewidth]{figures/ratio_comparison.png}
  \caption{Upper-to-lip motion ratio for VOCASET and the Disney Cinematic Corpus. The DCC ratio (1.21) is 8$\times$ higher than VOCASET (0.15).}
  \label{fig:ratio_comparison}
\end{figure}

Figure~\ref{fig:ratio_dist} shows the distribution of per-clip upper/lip ratios across the DCC, with the VOCASET reference marked as a dashed vertical line. All DCC clips fall well above the VOCASET baseline, with ratios ranging from approximately 0.8 to 1.6.

\begin{figure}[H]
  \centering
  \includegraphics[width=0.75\linewidth]{figures/ratio_dist.png}
  \caption{Distribution of upper/lip motion ratios across Disney Cinematic Corpus clips. The dashed line marks the VOCASET reference (0.15). All DCC clips exceed this baseline by a wide margin.}
  \label{fig:ratio_dist}
\end{figure}

% -------------------------------------------------------
% 6. DISCUSSION
% -------------------------------------------------------
\section{Discussion}

\paragraph{Limitations of the emotion classifier.}
The low human-model agreement (31.5\%) should not be interpreted as a failure of the classifier per se. The model was trained on general-domain text and achieves competitive performance on standard benchmarks. The low agreement reflects instead the inherent difficulty of cinematic speech as an affective signal: utterances such as \textit{``...Then leave.''} (Elsa, \textit{Frozen}) or \textit{``I wish you could have a door.''} (Antonio, \textit{Encanto}) carry emotional weight that is inseparable from the narrative context in which they occur. This motivates the framing of cinematic dialogue as a stress test rather than a standard evaluation benchmark.

\paragraph{FaceFormer amplifies the lower-face bias.}
Even on the neutral VOCASET sentences it was trained on, FaceFormer does not faithfully reproduce ground truth facial dynamics: it suppresses upper-face motion by 59\%, nearly twice the suppression applied to lip motion (35\%). This is not a global scaling effect, but it is a selective attenuation of upper-face dynamics that reflects the training objective. The standard evaluation metric, Lip Vertex Error (LVE), defined as the maximal L2 distance between predicted and ground truth lip vertices, penalizes lip error but places no constraint on upper-face expressivity; FaceFormer optimizes accordingly.

\paragraph{Cinematic dialogue demands a qualitatively different motion profile.}
The upper-to-lip ratio in the DCC (1.21) is 8$\times$ higher than in VOCASET (0.15). This gap is not a matter of degree but of kind: in neutral phonetic speech, lip movement dominates and upper-face movement is residual; in cinematic dialogue, the upper face is an expressive signal of comparable magnitude to the lips. A model trained exclusively on VOCASET has never been exposed to this regime.

\paragraph{Consistency across modules.}
Modules 1 and 2 tell a consistent story at different levels of analysis. The emotional gap between VOCASET and cinematic speech documented in Module 1 (a 35.6 percentage-point difference in neutral speech, a 31.5\% human-model agreement on cinematic text) maps directly onto the facial motion gap documented in Module 2: an 8$\times$ difference in upper/lip ratio, and a 59\% suppression of upper-face motion in the state-of-the-art system. The same structural bias that makes cinematic speech opaque to an emotion classifier makes it unexpressive in a facial animation system.

\paragraph{Limitations.}
The DCC is small ($n=73$) and curated rather than randomly sampled, which limits its generalizability as a benchmark. Module 2 analyses a single VOCASET test sequence; a full evaluation would require multiple sequences across subjects. The coordinate-based region definition, while topology-agnostic, is sensitive to the choice of coordinate thresholds and may include or exclude vertex boundary cases differently from anatomically grounded definitions. Future work should validate these findings on a broader set of sequences and with official LVE/FDD vertex indices.

% -------------------------------------------------------
% 7. CONCLUSION
% -------------------------------------------------------
\section{Conclusion}

This report has presented a two-module empirical investigation of cinematic speech as a stress test for facial animation systems. Module 1 established that VOCASET is structurally biased toward neutral speech: 68.5\% of deduplicated utterances are classified as neutral by an automated classifier, more than twice the proportion observed in the Disney Cinematic Corpus (32.9\%). The emotional content of cinematic dialogue is sufficiently complex to challenge state-of-the-art text classifiers, with human-model agreement of only 31.5\% on the DCC.

Module 2 extended this finding to the visual domain. FaceFormer, the dominant open-source system in speech-driven 3D facial animation, suppresses upper-face motion by 59\% elative to ground truth,  disproportionately more than the 35\% suppression applied to lip motion. The upper-to-lip motion ratio in cinematic dialogue (1.21) is 8$\times$ higher than in VOCASET (0.15), indicating that the facial motion demanded by cinematic speech is qualitatively different from what current systems are designed to produce.

Together, these results provide quantitative grounding for a hypothesis that has been acknowledged qualitatively in the literature but not yet measured: cinematic dialogue constitutes a genuine and largely unexplored stress test for speech-driven facial animation. Addressing the gap requires either training data that captures emotionally rich speech (as proposed by EmoVOCA~\cite{emoVOCA2025}) or explicit architectural conditioning on upper-face dynamics and emotional content.


% -------------------------------------------------------
% REFERENCES
% -------------------------------------------------------
\bibliographystyle{unsrt}
\begin{thebibliography}{99}

\bibitem{cudeiro2019capture}
D.\ Cudeiro, T.\ Bolkart, C.\ Laidlaw, A.\ Ranjan, and M.\ J.\ Black,
``Capture, learning, and synthesis of 3D speaking styles,''
in \textit{Proc.\ CVPR}, 2019.

\bibitem{fan2022faceformer}
E.\ Fan, F.\ Zhang, L.\ Feng, and G.\ Pillai,
``FaceFormer: Speech-driven 3D facial animation with transformers,''
in \textit{Proc.\ CVPR}, 2022.

\bibitem{sheng2023codetalker}
H.\ Sheng, Z.\ Mao, and L.\ Song,
``CodeTalker: Speech-driven 3D facial animation with discrete motion prior,''
in \textit{Proc.\ CVPR}, 2023.

\bibitem{witzig2025phonemenet}
S.\ Witzig et al.,
``PhonemeNet: Phoneme-level speech-driven facial animation,''
ETH Zurich, Computer Graphics Laboratory, 2025.

\bibitem{emoVOCA2025}
M.\ Daněček et al.,
``EmoVOCA: Speech-driven emotional 3D talking heads,''
\textit{arXiv preprint arXiv:2403.12886}, 2024.

\bibitem{hartmann2022}
J.\ Hartmann,
``Emotion English DistilRoBERTa-base,''
Hugging Face, 2022.
\url{https://huggingface.co/j-hartmann/emotion-english-distilroberta-base}

\bibitem{liu2019roberta}
Y.\ Liu et al.,
``RoBERTa: A robustly optimized BERT pretraining approach,''
\textit{arXiv preprint arXiv:1907.11692}, 2019.

\bibitem{fanelli2010realtime}
G.\ Fanelli, J.\ Gall, H.\ Romsdorfer, T.\ Weise, and L.\ Van Gool,
``A 3D audio-visual corpus of affective communication,''
\textit{IEEE Trans.\ Multimedia}, 12(6):591--598, 2010.

\bibitem{kim2021}
B.\ Kim and H.\ Bhatt,
``Implicit emotion detection in dialogue: Challenges and approaches,''
in \textit{Proc.\ EMNLP Findings}, 2021.

\bibitem{kaistami2025perceptual}
J.\ Kim et al.,
``Perceptually accurate 3D talking head generation: New definitions,
speech-mesh representation, and evaluation metrics,''
in \textit{Proc.\ CVPR}, 2025.

\end{thebibliography}

% -------------------------------------------------------
% APPENDIX
% -------------------------------------------------------

\appendix
\section{Fine-grained to Ekman Label Mapping}
\label{app:label_mapping}

Table~\ref{tab:label_mapping} reports the full mapping from fine-grained human annotation labels to the seven Ekman-aligned categories used by the emotion classifier, applied in the human-model agreement analysis of Section~\ref{sec:agreement}.

\begin{table}[H]
\centering
\small
\caption{Mapping from fine-grained DCC emotion labels to Ekman categories.}
\label{tab:label_mapping}
\begin{tabular}{p{10cm}l}
\toprule
\textbf{Fine-grained labels} & \textbf{Ekman} \\
\midrule
irony, deflection, cold, cold\_calculation, cold\_villainy, authority, wisdom, warmth, reassurance, innocence, control, protective, determination, empathy, protective\_love, denial, sarcasm
& neutral \\
\midrule
joy, excitement, tenderness, nostalgia, wonder, inspiration
& joy \\
\midrule
sadness, grief, melancholy, longing, regret, guilt, vulnerability, loneliness, hidden\_pain, numbness, self\_sacrifice, heartbreak, pleading, breakdown, inadequacy, exclusion, self\_doubt, suppressed\_sadness, emotional\_release
& sadness \\
\midrule
anger, defiance, frustration, repressed\_anger, injustice, confrontation
& anger \\
\midrule
fear, anxiety, panic, protective\_fear, desperation, repression, pressure
& fear \\
\midrule
disgust, manipulation, gaslighting, betrayal
& disgust \\
\midrule
surprise, shock, realization, inner\_conflict
& surprise \\
\bottomrule
\end{tabular}
\end{table}

\end{document}